\section{Conclusion}
Humanity need not choose between declaring alignment solved and concluding that nothing useful can be done. A model's intelligence, incentives, authority, and supporting evidence are distinct objects; legitimate authorization, adequate specifications, and reliable enforcement are distinct obligations. Keeping them separate exposes mistakes that reassurance and pessimism can both conceal.

This paper supplies a scoped assurance framework, conditional composition rules, an executable finite demonstration, and a local transactional prototype. Its strongest result is intentionally modest: a small, external, stateful boundary can preserve specified constraints across every modeled proposal sequence while still permitting useful work. The lifecycle ablations show why that boundary must preserve state across restarts and renewed permission. The specification control shows why even a useful, fully checked system can protect the wrong property. The new classifier study illustrates a second discipline: independence imposed by a simulation must not be promoted into independence of real safeguards. The next research task is to test transfer with independently challenged specifications, accountable grant procedures, realistic implementations, and adversaries strong enough to challenge the premises. Care-oriented training should be tested alongside that work, including whether it improves conduct when people refuse, disagree, or choose to leave. Successful control must protect affected people as well as serve the operator; otherwise it can make an unjust instruction more reliable.

Hope here is not a forecast. It is the existence of valuable work whose claims can be checked, whose mistakes can be corrected, and whose benefits do not depend on first resolving every question about machine minds. Intelligence can expand what we are able to propose. Evidence and legitimate human institutions should govern what we allow to take effect.
